\lecture{12}
\title{The EM Algorithm}

\begin{document}

\subsection{Computing the MLE}

In this lecture, we'll turn our attention to actually computing the MLE.
For very simple models, computing the MLE as a function of the sample data
$\{X_i\}_{i = 1}^n$ is very easy; recall
\href{/18.650/lectures/9/#defining-the-maximum-likelihood-estimator-mle}{Lecture 9}.

More generally, the MLE is the solution to an \emph{optimization} problem:
for what value $\theta$ is $\ell_n(\theta) := \sum_{i = 1}^n \log f_{\theta}(X_i)$ maximized?

There are two natural approaches to solving optimization problems.

\begin{quote}
    \textbf{Gradient Ascent.} Assuming $\ell_n'(\theta)$
    is computable, we can start at some guess $\theta_1$ and search for the maximum
    by picking some constant $\eta \in (0, 1)$ of our choice
    and updating our guesses $\{\theta_i\}_{i = 1}^{\infty}$ like so:
    \[ \theta_{i + 1} = \theta_i +
    \eta \cdot \ell_n'(\theta_i). \]
\end{quote}

\begin{quote}
    \textbf{Newton-Raphson (Root Finding).} The idea is to find a root of $\ell_n'(\theta)$
    and hope the critical point we find is a global maximum.
    The \href{https://en.wikipedia.org/wiki/Newton's_method}{Newton-Raphson algorithm}
    updates guesses $\{\theta_i\}_{i = 1}^{\infty}$ for the root like so:
    \[ \theta_{i + 1} = \theta_i - \frac{\ell'_n(\theta_i)}{\ell''_n(\theta_i)}. \]
\end{quote}

Fortunately, when our model is something like $\{\mathcal{N}(\theta, 1) \mid \theta \in \mathbb{R}\}$,
the log likelihood $\ell_n(\theta)$ is not very hard to optimize,
because its derivative is simple.
(This depends on normal distributions nicely having exponential PDFs!)

But there's a special kind of model
that is particularly hard to work with: \textbf{mixture models.}

\subsection{Mixture Models}

If we measure the heights of a random group of people,
a model like $\{\mathcal{N}(\theta, 1) \mid \theta \in \mathbb{R}\}$
won't work perfectly anymore.
One reason is that our sample is a \emph{mixture}
of men and women.  We might (simplistically) assume the following:
\begin{itemize}
    \item Any given sample has some probability $p$ of being from a woman,
    and $(1 - p)$ of being from a man.
    \item The height of a randomly chosen man
    has some PDF $f_{\mu_0}$ from the model $\{\mathcal{N}(\mu_0, 1) \mid \mu_0 \in \mathbb{R}\}$.
    \item The height of a randomly chosen woman
    has some PDF $f_{\mu_1}$ from the model $\{\mathcal{N}(\mu_1, 1) \mid \mu_1 \in \mathbb{R}\}$.
\end{itemize}
Our samples thus follow the PDF
$f = (1 - p)f_{\mu_0} + pf_{\mu_1}$.  This is a \textbf{mixture model},
with parameter $(\mu_0, \mu_1, p)$.  More generally\dots

\begin{quote}
    \textbf{Definition.} In a \textbf{mixture model},
    there are two models with parameters $\theta_0$ and $\theta_1$,
    along with corresponding PDFs $f_{\theta_0}$ and $g_{\theta_1}$.
    Every sample $i$ yields some data $(Y_i, Z_i)$, where:
    \[ Z_i \sim \mathrm{Ber}(p) ~~~ \text{ and } ~~~
    Y_i \sim \begin{cases}
    \mathbb{P}_{\theta_0} & \text{ if } Z_i = 0. \\
    \mathbb{P}_{\theta_1} & \text{ if } Z_i = 1.
    \end{cases} \]
    Note that we may only receive the data $\{Y_i\}_{i = 1}^n$
    and don't even know the value of $p$.
\end{quote}

There are two challenges with finding the MLE for mixture models.

\textbf{Challenge \#1.} We don't know the value of $p$,
so the log likelihood $\ell_n(\mu_0, \mu_1, p)$ has a floating $p$.
This means we have a third parameter to worry about now.

\textbf{Challenge \#2.} Even if we did assume, say, $p = \frac{1}{2}$,
the log likelihood still looks awful.
\[ \ell_n(\mu_0, \mu_1) = \sum_{i = 1}^n \log\left(
    \frac{1}{2} \left[\frac{1}{\sqrt{2\pi}}\exp\left(
        -\frac{1}{2}(Y_i - \mu_0)^2
    \right)\right] +
    \frac{1}{2} \left[\frac{1}{\sqrt{2\pi}}\exp\left(
        -\frac{1}{2}(Y_i - \mu_1)^2
    \right)\right]
\right) \]
The difficulty is that the logarithm of a sum does not simplify well at all.

\subsection{The EM Algorithm}

It turns out there is a strategy for computing the MLE of a mixture model:
the \textbf{EM Algorithm}.

\begin{quote}
    \textbf{Idea \#1.} If we know all of the $Z_i$,
    then estimating $(p, \theta_0, \theta_1)$ is doable.

    \emph{Reasoning:} Of course, if we know all the $Z_i$, we can just estimate $p$ directly.
    \[ \hat{p} = \frac{1}{n} \sum_{i = 1}^n Z_i. \]
    As for $(\theta_0, \theta_1)$,
    the rough idea is just to split the samples $\{Y_i\}_{i = 1}^n$
    into two halves based on their value of $Z_i$.
    Then just use one half to estimate $\theta_0$, and the other half to estimate $\theta_1$.
    Formally, this looks like:
    \[ \hat{\theta}_0 = \argmax_{\theta_0} \sum_{i = 1}^n (1 - Z_i) \log f_{\theta_0}(Y_i)
    ~~~ \text{ and } ~~~
    \hat{\theta}_1 = \argmax_{\theta_1} \sum_{i = 1}^n Z_i \log g_{\theta_1}(Y_i). \]
    And these are a lot easier to compute.
\end{quote}

Of course, we don't know all of the $Z_i$.
But maybe we can at least \emph{estimate} the $Z_i$\dots

\begin{quote}
    \textbf{Idea \#2.} If we have some prior estimate
    for $(p, \theta_0, \theta_1)$,
    then we can estimate $Z_i$ reasonably well.

    \emph{Reasoning:} Well, if we know $(p, \theta_0, \theta_1)$,
    then $Z_i$ can be estimated via Bayes' Rule.
    \[ \hat{Z}_i = \frac{g_{\theta_1}(Y_i) \cdot p}{
        g_{\theta_1}(Y_i) \cdot p + f_{\theta_0}(Y_i) \cdot (1 - p)
    }. \]
    Here, $\hat{Z}_i$ is not a discrete value in $\{0, 1\}$,
    but rather a continuous value in $[0, 1]$.
    Still, it's a good enough estimate that the reasoning in \textbf{Idea \#1} still holds.
\end{quote}

We have a problem though: \textbf{Idea \#1} relies on \textbf{Idea \#2},
and vice versa!  Here's the trick, though: we don't care.

\begin{quote}
    \textbf{EM Algorithm.} Start with some prior estimate
    $(\theta_0^{(0)}, \theta_1^{(0)}, p^{(0)})$.
    At time step $t$, make the following updates:
    \begin{itemize}
        \item \textbf{E Step.} In accordance with \textbf{Idea \#2},
        we estimate $\hat{Z}_i^{(t)}$ using last time step's estimates
        $(\theta_0^{(t - 1)}, \theta_1^{(t - 1)}, p^{(t - 1)})$
        and Bayes' Rule.
        \[ \hat{Z}_i^{(t)} = \frac{g_{\theta_1^{(t - 1)}}(Y_i) \cdot p^{(t - 1)}}{
            g_{\theta_1^{(t - 1)}}(Y_i) \cdot p^{(t - 1)} +
            f_{\theta_0^{(t - 1)}}(Y_i) \cdot (1 - p^{(t - 1)})
        }. \]
        \item \textbf{M Step.} In accordance with \textbf{Idea \#1},
        we compute $(\theta_0^{(t)}, \theta_1^{(t)}, p^{(t)})$
        using the estimates $\hat{Z}_i^{(t)}$ computed in the E Step.
        \[ p^{(t)} = \frac{1}{n} \sum_{i = 1}^n \hat{Z}_i^{(t)} ~~~ \text{ and } ~~~
        \theta_0^{(t)} = \argmax_{\theta_0} \sum_{i = 1}^n (1 - \hat{Z}_i^{(t)}) \log f_{\theta_0}(Y_i)
        ~~~ \text{ and } ~~~
        \theta_1^{(t)} = \argmax_{\theta_1} \sum_{i = 1}^n \hat{Z}_i^{(t)} \log g_{\theta_1}(Y_i). \]
    \end{itemize}
\end{quote}

\begin{remark}
    Suppose that instead of a mixture of just $2$ distributions,
    we had a mixture of $k$ distributions.

    There's an easy fix:
    instead of saying the $i^{\text{th}}$ sample only has the data $(Y_i, Z_i)$,
    say it has the data $(Y_i, Z_{i, 1}, Z_{i, 2}, \dots, Z_{i, k})$,
    where $(Z_{i, 1}, Z_{i, 2}, \dots, Z_{i, k}) \sim
    \mathrm{Categorical}(p_1, p_2, \dots, p_k)$ for all $i$,
    and $p_1 + p_2 + \dots + p_k = 1$.  Then:
    \begin{itemize}
        \item In the E step, make all estimates $\hat{Z}_{i, \ell}$
        using a slightly-more-involved Bayes' Rule.
        \item In the M step, update all probabilities $\{p_{\ell}\}_{\ell = 1}^k$
        and parameters $\{\theta_{\ell}\}_{\ell = 1}^k$
        by taking into account all estimates $\hat{Z}_{i, \ell}$.
    \end{itemize}
\end{remark}

It turns out the following is true.

\begin{theorem}{EM Algorithm Efficacy}
    The EM algorithm never decreases the log likelihood; that is,
    \[ \ell_n(\theta_0^{(t + 1)}, \theta_1^{(t + 1)}, p^{(t + 1)}) \geq
    \ell_n(\theta_0^{(t)}, \theta_1^{(t)}, p^{(t)}) ~ \text{ for all } t. \]
\end{theorem}

Proof omitted---it's fairly involved.

\end{document}